\chapter{大语言模型}

本章从模型机制的角度重新理解本书一直在使用的协作伙伴。理解 token、表示、注意力、训练与对齐，并不是为了从零训练大模型，而是为了更准确地判断提示、上下文、工具调用和模型输出之间的关系。

\section{大语言模型的技术基础}
\label{sec:llm_foundation}

\subsection{为什么在这里讲大语言模型}

在前文中，我们已经依次介绍了前馈神经网络、循环神经网络、注意力机制、Transformer、卷积神经网络和生成模型。到了这一阶段，再引入大语言模型（Large Language Models, LLMs）是很自然的，因为它并不是一类凭空出现的新技术，而是在 Transformer 架构、海量数据预训练和大规模工程实践基础上发展起来的。

如果从技术脉络来看，卷积网络强调局部结构建模，图神经网络强调关系传播，注意力机制强调动态加权，而大语言模型则是在这些思想之上，进一步通过 Transformer 的统一序列建模能力和大规模预训练，把“理解”和“生成”两类任务尽可能合并到同一个框架中。

需要特别说明的是，Transformer 这一核心架构由 Vaswani 等人在 2017 年论文《Attention Is All You Need》中提出。该工作首次系统地用纯注意力机制取代了传统的循环结构和卷积结构，并给出了 encoder-decoder 形式的标准 Transformer。此后，BERT 将 Transformer 推向了双向表示学习，GPT 路线则将 decoder-only 架构与大规模自回归预训练结合起来，最终推动了今天大语言模型的发展。

因此，本节的目标不是把大语言模型说成一个“神秘黑箱”，而是帮助读者从已经掌握的 Transformer 与注意力知识出发，理解它为什么能工作、如何训练、如何调用，以及它在工程场景中能做什么、不能做什么。Q/K/V、多头注意力和自回归生成的细节已在前一章展开，本章只在必要处简要回顾，并把重点放在大语言模型的训练范式、应用方式和工程风险上。

\subsection{什么是大语言模型}

所谓大语言模型，通常指的是参数规模大、训练数据规模大、训练计算量大，并且具备较强迁移能力的一类语言模型。它的任务不再局限于某一个单独的分类或回归问题，而是希望在统一的文本建模框架下，同时支持问答、摘要、翻译、改写、信息抽取、代码生成、对话和推理等多种任务。

从更广义的角度看，大语言模型也可以被视为基础模型（Foundation Model）在文本领域的一种典型形式。其核心思想是：先在海量通用语料上进行预训练，学习尽可能丰富的语言模式和知识表示，再通过监督微调、提示词、对齐训练或检索增强等方式，把模型迁移到具体应用场景中。

如果把传统任务模型和大语言模型做一个直观对比，那么可以这样理解：

\begin{itemize}
    \item 传统任务模型更像“为一道题专门练一个模型”；
    \item 大语言模型更像“先学通用语言能力，再去做很多题”。
\end{itemize}

这种差异意味着：大语言模型更强调通用性、迁移性和生成能力，而不是只针对某一个封闭任务训练一个专门网络。

\subsection{大语言模型的最小输入单位：Token、Embedding 与位置编码}

为了让初学者也能明白“大模型的输入到底是什么”，我们先从最基础的输入单位开始。

对于大语言模型来说，原始输入文本并不是直接以“句子”的形式进入模型，而是要先被切分成若干个基本单元，通常称为 \textbf{token}。一个 token 可以是一个词、一个子词、一个汉字片段、一个标点，甚至是某种特殊符号。不同模型使用的 tokenizer 规则可能不同，但核心思想都是一样的：\textbf{先把文本离散化为 token 序列，再把 token 映射为向量。}

假设输入句子被切成长度为 $L$ 的 token 序列：
\[
x = (x_1, x_2, \dots, x_L)
\]
其中每个 $x_t$ 都是一个 token 的编号。接着，模型会用嵌入矩阵把它们映射为连续向量：
\[
\mathbf{e}_t = \mathbf{E}[x_t]
\]
如果模型维度为 $d_{\text{model}}$，那么每个 token 对应的向量维度就是 $d_{\text{model}}$。

因此，整句输入最终会变成一个矩阵：
\[
X \in \mathbb{R}^{L \times d_{\text{model}}}
\]

这一步非常关键，因为它解释了 Transformer 的真实输入到底是什么：\textbf{不是一句自然语言本身，而是一个“序列长度 $\times$ 向量维度”的矩阵。}换句话说，只要知道一句话有多少个 token、每个 token 是多少维向量，就能写出输入矩阵的形状。

\begin{figure}[H]
  \centering
  \resizebox{\textwidth}{!}{\begin{tikzpicture}[font=\small, >=Stealth, t/.style={draw, rounded corners, fill=yellow!20, minimum height=6mm}, tid/.style={draw, fill=blue!10, minimum width=9mm, minimum height=6mm}]
\node[anchor=west] at (-0.3,0.9) {原始文本：“早高峰匝道排队回溢”};
\foreach \tok/\i/\x in {早/1024/0, 高峰/5391/1.4, 匝道/20877/2.8, 排队/9142/4.2, 回溢/31006/5.6} {
  \node[t] (t\i) at (\x,0) {\tok};
  \node[tid] (d\i) at (\x,-1.1) {\scriptsize \i};
  \draw[->] (t\i) -- (d\i);
  \foreach \k in {0,...,3} \fill[green!\the\numexpr 15+15*\k\relax] (\x-0.35+0.175*\k,-2.4) rectangle ++(0.175,0.6);
  \draw (\x-0.35,-2.4) rectangle ++(0.7,0.6);
  \draw[->] (d\i) -- (\x,-1.8);
}
\node[anchor=east, font=\scriptsize] at (-0.6,0) {分词 (token)}; \node[anchor=east, font=\scriptsize] at (-0.6,-1.1) {token 编号}; \node[anchor=east, font=\scriptsize] at (-0.6,-2.1) {嵌入向量 $+$ 位置编码};
\draw[decorate, decoration={brace, amplitude=5pt, mirror}] (-0.4,-2.6) -- node[below=6pt, font=\scriptsize] {输入矩阵 $X\in\mathbb{R}^{L\times d_{\text{model}}}$：$L$ 个 token，每个是 $d_{\text{model}}$ 维向量（编号仅为示意）} (6.0,-2.6);
\end{tikzpicture}
}
  \caption{从文本到 Transformer 输入矩阵：分词、编号与向量化}
  \label{fig:ch11_orig1}
\end{figure}

图~\ref{fig:ch11_orig1} 应重点看两个转换：第一步是从文字到 token 编号，第二步是从 token 编号到连续向量。Transformer 后续处理的不是原始文字，而是由这些向量按顺序堆叠成的矩阵；因此，序列长度决定矩阵有多少行，模型维度决定每一行有多少列。

然而，Transformer 本身不像 RNN 那样自带顺序递归结构，所以它还需要额外加入位置信息。最经典的做法是把位置编码与 token embedding 相加：
\[
\mathbf{z}_t = \mathbf{e}_t + \mathbf{p}_t
\]
这样模型既知道“这个位置是什么 token”，也知道“这个 token 出现在序列中的什么位置”。Transformer 论文中给出了经典的正弦/余弦位置编码形式。

\subsection{Transformer 机制回顾：本章只保留大模型视角}

上一章已经详细介绍了 Q/K/V 的来源、缩放点积注意力、多头注意力、三类 Transformer 结构，以及自回归生成的逐 token 循环。本章不再重复这些推导，只保留与大语言模型直接相关的几条结论：

\begin{itemize}
    \item 文本先被切分为 token，再映射为向量矩阵 \(X\in\mathbb{R}^{L\times d_{\text{model}}}\)；
    \item self-attention 让序列中不同 token 直接建立内容相关性；
    \item 多头注意力让模型在多个表示子空间中并行建模关系；
    \item decoder-only 结构适合“根据已有上下文预测下一个 token”的生成任务；
    \item 因果遮罩保证模型在生成第 \(t\) 个位置时不能偷看未来内容。
\end{itemize}

因此，读者在本章阅读“预训练”“指令微调”“上下文学习”“Agent 工具调用”等概念时，可以始终把它们放回同一条主线中理解：大语言模型的底层仍然是 Transformer 序列模型，只是经过了更大规模的数据、参数和工程系统扩展。

\subsection{大语言模型的训练范式：预训练、监督微调与对齐}

大语言模型之所以能够表现出强大的迁移能力，不只是因为参数变大了，更重要的是训练范式变了。与传统“一个任务一个模型”的思路不同，大语言模型通常采用三阶段范式：

\paragraph{预训练}

预训练阶段通常在海量通用语料上进行。对于 decoder-only 语言模型，最常见的目标是“下一个 token 预测”：
\[
\mathcal{L}_{LM}
=
- \sum_{t=1}^{T}\log p(x_t \mid x_{<t})
\]

这种训练目标看起来简单，但当数据量和模型规模足够大时，模型会在不断预测下一个 token 的过程中，逐渐吸收语言结构、常识知识、代码模式以及任务形式。GPT 系列正是这一路线的典型代表，而 GPT-3 则系统展示了规模扩展带来的 few-shot / in-context learning 能力。

\paragraph{监督微调（SFT）}

完成预训练后，模型虽然已经具备一般语言建模能力，但未必能很好地“听指令”。因此，常常还要在人工构造的指令—回答数据上进行监督微调：
\[
\mathcal{L}_{SFT}
=
- \sum_t \log p(y_t \mid x, y_{<t})
\]

通过这一阶段，模型能够更好地学会摘要、翻译、解释、抽取、代码补全等任务的响应风格。

\paragraph{对齐训练（Alignment）}

即使经过 SFT，模型输出仍然可能不符合人类偏好，或者出现不安全、无帮助、事实不可靠的回答。因此，还需要进一步进行对齐训练，让模型输出更符合人类偏好与安全要求。常见思路包括 RLHF 和 DPO 等。

\paragraph{参数高效微调}

当模型变得很大时，直接微调整个模型成本很高，因此工程上经常使用 LoRA、Adapter、Prompt Tuning、Prefix Tuning 等参数高效微调方法，让模型只调整一小部分新增参数就能适配具体领域。

\subsection{为什么大语言模型有效}

大语言模型有效，通常可以从四个角度来理解。

第一，\textbf{规模效应}。随着参数规模、训练数据规模和训练计算量同步增大，模型能够学习到更复杂的统计规律和更丰富的知识表示。GPT-3 论文就是这一点的代表性例子。

第二，\textbf{统一建模方式}。许多任务都可以统一成“根据上下文预测后续 token”或者“根据指令生成目标输出”的形式，这降低了任务之间的结构壁垒。

第三，\textbf{上下文学习（in-context learning）}。模型即使不更新参数，也可能根据 prompt 里的任务描述和示例临时适应一个新任务，这也是大语言模型非常独特的一点。

第四，\textbf{预训练带来的通用先验}。模型在下游任务中表现出来的能力，并不只来自下游小数据集，而是建立在大规模先验学习基础上的。

\subsection{大语言模型与前文知识点的关系}

从本书整体结构看，大语言模型并不是与前文割裂的主题。

它和 CNN 的关系在于：CNN 强调局部结构与层次特征提取，而大模型中的视觉分支仍然常常继承这些思想，或者通过 Vision Transformer 把图像切成 patch token 再送入 Transformer。

它和 GCN / GAT 的关系在于：GCN/GAT 强调图上的邻域聚合，而 self-attention 则在全局范围内根据内容相关性动态建立连接。一个依赖显式图结构，一个依赖动态学习到的相关性，但“信息聚合”的核心思想是一致的。

它和前文 attention 的关系就更直接了：如果说前文介绍的是 attention 的机制，那么大语言模型本质上就是把这一机制规模化、层次化、工程化之后形成的统一框架。

\section{应用、局限与协作系统}

\subsection{多模态大模型与工程意义}

随着研究发展，大模型已经不局限于文本。多模态大模型试图统一处理文本、图像、语音、视频、表格乃至图结构等多种输入。例如，一个工程智能系统可以同时读取巡检图片、传感器时序、结构拓扑和运维文本，然后输出风险判断或维护建议。

对于土木工程与交通工程，这种能力尤其重要，因为真实工程系统本身就是多源异构的：结构健康监测既有时序传感器数据，也有巡检图像和结构关系；智能交通既有时空流量图，也有监控视频和事件文本。大模型的价值不一定是“完全替代专用模型”，而更可能是作为统一的知识整合与交互层，把 CNN、GCN、时序模型、文档资料和外部工具连接起来。

\begin{figure}[H]
  \centering
  \resizebox{\textwidth}{!}{\begin{tikzpicture}[font=\scriptsize, >=Stealth, d/.style={draw, rounded corners, fill=gray!10, minimum width=24mm, minimum height=7mm}, e/.style={draw, rounded corners, fill=blue!12, minimum width=22mm, minimum height=7mm}]
\foreach \n/\lab/\enc/\y in {img/巡检照片、监控视频/视觉编码器（CNN、ViT）/1.8, ts/传感器与交通流时序/时序编码器/0.6, g/路网与结构拓扑/图编码器（GCN）/-0.6, txt/维护记录、事件文本/文本编码器（Transformer）/-1.8} {
  \node[d] (\n) at (0,\y) {\lab}; \node[e] (e\n) at (3.6,\y) {\enc}; \draw[->] (\n) -- (e\n);
}
\node[draw, rounded corners, fill=orange!20, minimum height=40mm, minimum width=18mm, align=center] (fuse) at (7.2,0) {共享表示\\空间与\\跨模态融合};
\foreach \n in {img,ts,g,txt} \draw[->] (e\n) -- (fuse);
\node[draw, rounded corners, fill=green!20, align=center] (out) at (10.0,0) {风险判断\\维护建议\\问答与报告};
\draw[->] (fuse) -- (out);
\end{tikzpicture}
}
  \caption{多模态模型处理工程数据的一般结构}
  \label{fig:ch11_orig2}
\end{figure}

图~\ref{fig:ch11_orig2} 强调的是工程数据并非只有文本一种形态。图像可以来自巡检照片，时序可以来自传感器或交通流量，图结构可以表示路网或结构连接关系，文本则可能是维护记录或事件报告。多模态模型的核心任务，是先把这些异构数据编码到可比较的表示空间，再在融合层中完成跨模态信息整合。

\subsection{大语言模型的局限性与风险}

尽管大语言模型能力很强，但它并不是万能的。在工程场景里尤其要强调以下几点：

\begin{itemize}
    \item 幻觉：模型可能生成语言流畅但事实错误的内容；
    \item 时效性不足：模型知识主要来自训练语料，不一定覆盖最新事实；
    \item 计算成本高：训练和部署都比较耗算力与存储；
    \item 可解释性有限：内部决策过程复杂，不易完全解释；
    \item 工程可靠性要求高：在土木、交通等高风险场景中，不能不加约束地把模型输出当最终决策。
\end{itemize}

因此，在实际工程应用中，大语言模型更适合作为\textbf{辅助决策工具}，并与 RAG、规则库、知识图谱、数据库、仿真模型和人工审核协同使用。

\subsection{AI 时代为什么仍然要学习代码}

进入大模型时代后，很多代码已经可以由 AI 辅助生成。许多 AI 编程工具都可以根据自然语言描述生成函数、解释报错、补全测试，甚至辅助修改一个项目中的多个文件。这会带来一个很自然的问题：既然 AI 已经能写代码，为什么还要学习本书中的机器学习与 AI 系统知识？

答案是：\textbf{AI 可以帮助生成代码，但不能替代人理解问题、校验结果和判断代码是否可信。} 如果读者不理解模型、数据、目标函数和评价指标，那么即使 AI 生成了一段能运行的程序，也很难判断它是否真的解决了正确的问题。在医疗、金融、交通与公共服务等高影响场景中，错误代码可能进一步导致错误预测与错误决策，因此更需要人具备基本的代码阅读和验证能力。

因此，本书中的代码实训有两层目的。第一层，是提供规范、可运行、可复现实验结果的示例，让读者真正理解数据如何进入模型、模型如何训练、结果如何评价。第二层，是帮助读者学会与 AI 协作：可以让 AI 根据任务描述生成初稿，也可以让 AI 解释代码、补充注释、生成测试，但最终仍要由读者检查数据来源、模型逻辑、输出指标和可视化结果是否合理。

在前几个实训中，可以尝试使用较完整的提示词。例如：

\begin{quote}
请用 Python 写一个可直接运行的机器学习示例，使用 scikit-learn 自带的真实数据集，完成数据划分、标准化、模型训练、测试集评价和结果可视化。代码需要包含清晰注释，并解释每个评价指标的含义。
\end{quote}

当读者熟悉这种模式后，后续练习可以逐步减少提示词细节，只给出任务目标，让读者自己判断 AI 生成的代码是否完整、是否使用了可靠数据、是否存在数据泄露、是否正确解释了结果。这样，AI 不再只是“代写代码”的工具，而是帮助读者练习建模思维、调试能力和工程判断的助手。

\subsection{Agent、MCP、Skills 与 ReAct：不要被新名词吓住}

近年来，大模型外层出现了很多新框架和新术语，例如 AI Agent、工具调用、RAG、MCP、Skills、Workflow、ReAct 等。这些概念看起来很复杂，但如果回到本质，大部分系统仍然建立在 Transformer 大语言模型和“下一个 token 预测”之上。区别在于：模型外面增加了提示词、工具接口、记忆、检索、执行反馈和安全约束，使它能够完成更长流程的任务。

可以先用一个生活化类比理解：普通聊天模型像一个只能在纸面上回答问题的助手；Agent 则像一个带着工具箱的助手。前者只能根据已有上下文回答，后者还可以查资料、打开文件、运行代码、把结果拿回来继续判断。大模型本身仍然负责“读上下文并生成下一步文字”，但系统把某些文字解释成工具调用，于是模型就能和外部世界发生交互。

可以把一个典型 AI Agent 理解为以下循环：

\begin{enumerate}
    \item 用户提出任务，例如“读取数据、训练模型并生成报告”；
    \item 系统提示词规定模型的角色、权限、输出格式和安全边界；
    \item 模型根据上下文生成下一步计划或工具调用；
    \item 外部程序执行工具，例如读取文件、运行代码、查询数据库或调用 API；
    \item 工具结果返回给模型，模型再判断下一步；
    \item 重复上述过程，直到任务完成或需要人工确认。
\end{enumerate}

ReAct 框架可以理解为“推理（Reasoning）+行动（Acting）”的交替：模型先根据当前信息思考下一步，再调用工具执行动作，然后根据反馈继续推理。MCP（Model Context Protocol）一类协议，则更像是为模型和外部工具之间建立统一接口，让模型能够以标准方式访问文件、数据库、网页、代码仓库或专业软件。Skills 则可以理解为封装好的任务说明和操作流程，告诉模型在特定任务中应该按什么步骤工作。

把这些术语拆开看，就没有那么神秘：
\begin{itemize}
    \item \textbf{Agent}：不是一个新模型，而是“大模型 + 工具 + 执行循环”的系统；
    \item \textbf{ReAct}：不是魔法推理，而是“先想一下该做什么，再行动，再根据反馈继续想”；
    \item \textbf{MCP}：不是新的神经网络结构，而是让模型按统一方式连接外部工具和数据源的接口协议；
    \item \textbf{Skills}：不是模型突然学会了新知识，而是把某类任务的步骤、约束和经验写成可复用说明，让模型更稳定地照着做。
\end{itemize}

例如，让 Agent 帮忙“检查一段 CNN 代码是否正确”时，它可能先阅读代码，发现需要知道张量形状；然后运行脚本，读取报错或输出；再根据输出判断卷积层通道数是否匹配；最后给出修改建议。这个过程中，大模型并没有脱离 next token prediction，只是每一步生成的内容会影响下一次工具调用和下一轮上下文。

这些框架确实提高了大模型的工程能力，但它们并没有改变最底层的学习范式：模型仍然是根据上下文预测下一个 token，只是上下文中多了系统提示词、工具返回结果、历史步骤和约束规则。理解这一点后，读者就不会被新名词吓住。无论外层框架如何变化，判断一个大模型系统是否可靠，仍然要回到几个基本问题：输入是什么，模型能看到哪些上下文，工具是否真实执行，结果是否可验证，错误是否会被及时发现。

\section{案例：用大语言模型对交通投诉文本进行分类与信息抽取}
\label{sec:ch11_case}

本案例与贯穿案例 C 相关：客流预测回答的是“有多少人”，而乘客反馈回答的是“乘客遇到了什么问题”。MTA 在纽约州开放数据平台上公开了自 2018 年起的客户反馈记录，每条记录包含类别、子类别、线路、地点和时间等结构化字段，但出于隐私保护，公开版本不包含反馈的原文。这是公共交通数据的常见情形：真正需要大语言模型处理的自由文本，往往只存在于机构内部。

为了在公开数据上完整演示这一方法，本案例使用另一个真实的公开文本数据源：美国国家公路交通安全管理局（NHTSA）的车辆安全投诉数据库。每条投诉包含车主撰写的问题描述原文，以及数据库中记录的故障部件类别（如制动系统、转向系统、安全气囊等），可以作为检验分类结果的参照标签。我们通过 NHTSA 的公开接口下载了 16 款常见车型 2018—2021 年款的全部投诉，共 29,492 条；保留只涉及一个部件、且属于 8 个常见部件类别的 10,571 条，从中按类别分层随机抽取 160 条（每类 20 条）作为评价集，其余 10,411 条作为训练传统基线的样本。数据来自 NHTSA 的公开投诉接口\cite{nhtsa2026complaints}。配套代码见 \texttt{code/misc/fetch\_nhtsa.py}、\texttt{make\_nhtsa\_gold.py} 与 \texttt{ch11\_nhtsa\_eval.py}；由于数据库持续新增投诉，书中使用的 160 条评价集和模型输出随代码一并保存在 \texttt{code/misc/nhtsa\_eval/} 中。

\subsection{第一步：先定义“正确答案”}

在编写任何 Prompt 之前，应当先与业务人员一起制定分类说明：列出类别（例如，对于地铁乘客反馈，可以是站内拥挤、设备故障、安检排队、票务问题、服务态度、环境卫生、列车延误等），并为每一类写明定义、典型例子和容易混淆的边界情况。例如，“安检口排了二十分钟”属于安检排队而不是站内拥挤；“换乘通道人太多走不动”属于站内拥挤。

随后，由两名熟悉业务的人员独立标注一批随机抽取的文本（例如 300 至 500 条），计算两人的一致率和 Cohen's $\kappa$，并通过讨论确定有分歧样本的最终标签，形成“金标准”评价集。这一步的意义在于：它给出了“人与人之间能达到多高的一致性”这一参照。如果两名熟练人员本身就有 10\% 的分歧，就不应期待模型达到 99\% 的准确率，也不应把模型与金标准的每一处差异都视为模型的错误。对 NHTSA 数据而言，官方归类的部件字段可以作为参照，但它本身也是人工归类的结果，同样需要抽样核对。

\subsection{第二步：设计可验证的输出}

\begin{quote}
\textbf{角色与任务：}你是交通投诉文本的分类助手。请阅读下面的文本，按照分类说明完成分类，并抽取关键信息。

\textbf{分类说明：}（此处附完整的类别定义、典型例子和边界规则，以及若干条覆盖易混淆类别的示例及其正确输出。）

\textbf{抽取规则：}
\begin{enumerate}
  \item \texttt{location}：只有当文本中明确出现车站或地点名称时才填写，并且必须是附件名录中的标准名称之一；文本中未出现或无法唯一对应时，填写 \texttt{null}，不得根据线路或地标推测；
  \item \texttt{time\_period}：只能从“早高峰、平峰、晚高峰、夜间、未提及”中选择，不得生成具体时刻；
  \item \texttt{category}：从给定类别中选择 1 个主类别；确实涉及两类时，可在 \texttt{secondary} 中填写第二类；
  \item \texttt{evidence}：摘录原文中支持分类的原句，不得改写。
\end{enumerate}

\textbf{输出格式：}只输出一个 JSON 对象，字段为 \texttt{category}、\texttt{secondary}、\texttt{location}、\texttt{time\_period}、\texttt{evidence}，不要输出任何其他文字。
\end{quote}

这份 Prompt 中有三处关键设计。其一，\texttt{evidence} 字段要求模型摘录原文，使每个分类决定都可以被人快速核对，并且可以用程序检查摘录是否真的出现在原文中。其二，对地点和时间明确规定了“不知道就填 \texttt{null}”，并限定了取值范围。这直接针对大语言模型最典型的错误：在信息不足时生成一个看起来合理的具体值——例如乘客只说“体育馆那站”，模型却填写一个听起来合理、但名录中并不存在的站名。其三，输出为固定结构的 JSON，可以被程序逐字段校验。

在正式运行之前，还有一件不属于建模、但必须完成的事：投诉文本中常含有姓名、电话、车牌或卡号。所有文本在送入模型之前应先脱敏；如果使用外部模型服务，还需要确认数据使用条款和所在机构的数据安全要求。

\subsection{第三步：在金标准上评价}

评价至少应包括以下几项：与低成本的传统基线比较；分类别报告 Precision 与 Recall 以及宏平均 F1，而不只是总体准确率；核对抽取字段（地点是否在名录中、原文未提及时是否填写了 \texttt{null}、\texttt{evidence} 是否确实是原文的子串）；对同一批文本重复运行，检查输出是否稳定；以及输出格式的合规率。

在 NHTSA 评价集上，我们按上述方法做了一次实际评价。大语言模型部分由 Claude 在本书写作过程中完成：先固定分类说明（8 个类别的定义与 5 条边界规则，见 \texttt{code/misc/nhtsa\_eval/nhtsa\_guide.txt}），再逐条阅读不含标签的 160 条投诉并给出类别与证据摘录，全部完成之后才与参照标签比对。传统基线为词级 TF-IDF 加逻辑回归，分别用 500 条、2,000 条和全部 10,411 条训练样本训练。结果如表~\ref{tab:ch11_eval} 所示。

\begin{table}[H]
\centering
\caption{NHTSA 投诉部件分类的评价结果（160 条评价集，8 类各 20 条）}
\label{tab:ch11_eval}
\begin{tabular}{lcc}
\toprule
方法 & 准确率 & 宏平均 F1 \\
\midrule
TF-IDF + 逻辑回归（500 条训练样本） & 0.750 & 0.749 \\
TF-IDF + 逻辑回归（2,000 条训练样本） & 0.806 & 0.806 \\
TF-IDF + 逻辑回归（10,411 条训练样本） & 0.875 & 0.875 \\
大语言模型（分类说明 + 边界规则，无训练样本） & 0.925 & 0.925 \\
\bottomrule
\end{tabular}
\end{table}

大语言模型在没有使用任何训练样本的情况下，超过了用一万余条标注样本训练的传统基线；160 条证据摘录全部是原文的逐字子串。但逐条审阅 12 条不一致的样本，比准确率本身更有价值：

\begin{itemize}
  \item \textbf{我们写的边界规则与参照标签的口径不一致。}分类说明中规定“熄火或失去动力，除非明确归因于变速器或电气，否则归为发动机”。有 3 条“熄火”投诉因此被归为发动机，而参照标签是动力传动系统。进一步统计训练集发现，提到熄火（stall）的 472 条投诉中，参照标签为发动机的只占 68\%，动力传动系统占 18\%，电气系统占 12\%。换言之，这类投诉在参照标签中本身就没有统一口径。
  \item \textbf{参照标签并非绝对正确。}一条“头枕无故弹出”的投诉，参照标签为安全气囊；一条“安全气囊和安全带都没有起作用”的碰撞投诉，参照标签为安全带。它们更像是投诉在归类时的取舍，而不是模型“答错了”。提到自动紧急制动或“幻影刹车”的 81 条训练集投诉中，也有 12\% 没有被归为制动系统。
  \item \textbf{真正的模型错误集中在多部件描述上。}例如一条同时描述“踩刹车却加速”和“气囊未弹出”的投诉，模型按“造成安全风险的部件”选择了气囊，而参照标签选择了制动系统。这类情况需要在分类说明中补充更具体的优先级规则。
\end{itemize}

这次评价没有报告稳定性：同一会话中的“重复分类”并不是独立的重复运行。读者可以用 \texttt{ch11\_llm\_text\_classification.py} 通过 API 对同一评价集重复运行多次，统计类别发生变化的样本。还需要强调的是，本次评价由参与本书写作的模型完成，模型、版本和提示方式不同，结果也会不同；它展示的是评价方法，而不是任何模型的性能保证。

\subsection{第四步：确定 AI 在流程中的角色}

根据评价结果，一个稳妥的流程通常是：大语言模型作为执行者完成全部文本的初步分类与抽取；程序化规则检查格式、名录匹配和证据摘录；名录无法匹配、多次运行不一致或类别为“其他”的样本进入人工复核队列；此外定期从其余样本中随机抽检，持续监测准确率是否随时间下降（例如出现了新的投诉类型）。

结构化之后的反馈数据，才能与其他数据相互印证。以地铁为例，把“站内拥挤”类反馈按车站和时段汇总后，可以与第 16 章的客流预测结果对照：如果某座车站的拥挤反馈明显多于其客流规模所对应的水平，一个值得现场核实的假设是该站存在出入口或换乘通道等设施瓶颈——拥挤感未必来自客流总量。这类假设只能作为调查线索，而不能直接作为结论。

\subsection{本案例中人的判断}

\begin{itemize}
  \item 与一线人员共同制定分类说明，并用两人独立标注建立金标准和人际一致性参照；
  \item 在 Prompt 中限定取值范围、要求摘录证据，并对“不知道”给出明确出口；
  \item 在数据送入模型之前完成脱敏，并确认模型服务的使用方式符合数据安全要求；
  \item 与低成本的传统基线比较，并逐条审阅错误样本，而不是只看准确率；
  \item 用程序化规则和人工抽检构建持续验证机制，把大语言模型放在“执行者 + 人工复核”的位置上；
  \item 把文本分析的结果作为线索与其他数据相互印证，而不是把投诉数量直接等同于问题的严重程度。
\end{itemize}

大语言模型可以把过去只能抽样阅读的文本变成可以全量分析的结构化数据。但这种价值的前提，是人先定义了什么是正确答案，并设计了发现错误的机制。

% ----------------------------------------------------------------
\section{本章小结}
% ----------------------------------------------------------------

\begin{itemize}
  \item 大语言模型以 Token、Embedding、位置编码和 Transformer 为基础，经过预训练、监督微调与对齐获得通用能力。
  \item 它的局限包括幻觉、时效性、成本与可解释性；在工程场景中更适合作为与检索、规则、数据库和人工审核协同的辅助工具。
  \item Agent、MCP、Skills 与 ReAct 等框架没有改变“根据上下文预测下一个 token”的本质，判断系统是否可靠仍要回到输入、上下文、工具执行与验证。
  \item 用大语言模型处理文本时，应先定义金标准，设计可核查的输出（限定取值、摘录证据、允许 null），并与低成本基线比较、建立持续抽检机制。
\end{itemize}

% ----------------------------------------------------------------
\section{实战练习}
% ----------------------------------------------------------------

\subsection*{练习一}

选取 200 条公开的交通投诉文本（例如 NHTSA 车辆安全投诉），由两人独立标注并计算一致性，建立一个小型金标准。

\subsection*{练习二}

使用 \texttt{code/misc/ch11\_llm\_text\_classification.py}，比较大语言模型与 TF-IDF 基线的宏平均 F1，并统计 evidence 字段不是原文子串的比例。

\subsection*{练习三}

对同一批文本重复运行 3 次，找出类别不稳定的样本，并与人工标注存在分歧的样本进行对照。

